\documentclass{book}
\usepackage[utf8]{inputenc}
\usepackage{amsmath}
\usepackage{amsfonts}
\usepackage{esint}
\usepackage{geometry}
\usepackage{color}
\usepackage{fancyhdr}
\usepackage{ctable}
\usepackage{fancybox}
\usepackage{tabularx}
\usepackage{array}
\usepackage{booktabs}
\usepackage{subcaption}
\usepackage[french]{babel}
\usepackage{dsfont}
\usepackage{setspace}
\usepackage[french]{minitoc}
\usepackage{multicol}
\usepackage{multirow}
\usepackage[hidelinks]{hyperref}
\usepackage{graphicx}
\usepackage[T1]{fontenc}
\usepackage{xcolor}
\usepackage{listings}

\geometry{top=2.3cm, bottom=2.3cm, left=2.5cm, right=2.5cm}

% \addtocounter{tocdepth}{3}
% \setcounter{secnumdepth}{3}


\definecolor{codegreen}{rgb}{0,0.6,0}
\definecolor{codegray}{rgb}{0.5,0.5,0.5}
\definecolor{codepurple}{rgb}{0.58,0,0.82}
\definecolor{backcolour}{rgb}{0.95,0.95,0.92}

\lstdefinestyle{mystyle}{
  backgroundcolor=\color{white}, commentstyle=\color{codegreen},
  keywordstyle=\color{magenta},
  numberstyle=\tiny\color{codegray},
  stringstyle=\color{codepurple},
  basicstyle=\ttfamily\footnotesize,
  breakatwhitespace=false,         
  breaklines=true,                 
  captionpos=b,                    
  keepspaces=true,                 
  numbers=left,                    
  numbersep=5pt,                  
  showspaces=false,                
  showstringspaces=false,
  showtabs=false,                  
  tabsize=2
}

\lstset{style=mystyle}

\begin{document}

%%%%%%%%%%%%%%%%%%%%%%%%%%%%%%%%%%%%%%%%%%%%%%%%%%%%%%
%%%%%%%%%%%%%%%%%%%% PRÉSENTATION %%%%%%%%%%%%%%%%%%%%
%%%%%%%%%%%%%%%%%%%%%%%%%%%%%%%%%%%%%%%%%%%%%%%%%%%%%%

\begin{titlepage}

    \unitlength 1cm
    \begin{center}
    
    \vspace*{1cm}

    \includegraphics[scale=0.6]{figures/logo_ico.png}
    
    \vspace{2cm}
    
               {\Large Diplôme de Qualification en Physique Radiologique et Médicale\\}
               
    \vspace{2cm}           
    
    
    \rule{16cm}{0.7pt}
    
    \vspace{12pt}
               
               {\LARGE \bf Faisceaux de photons de haute énergie : cas des petits faisceaux\\}
               
    \vspace{12pt}
    \rule{16cm}{0.7pt}

    \vspace{2cm}

                {\large Fiche n°3b}
    
    \vspace{1.5cm}

               {\Large\bf {Alexandre \textsc{Rintaud}}}
    
    \vspace{1.5cm}
    
    \end{center}
    
    Encadrants :
    
    \small {
    \begin{tabular}{llr}\\
    \textbf{Stéphanie \textsc{Josset}}   &  &  \\
      Physicienne médicale, \textsc{Centre René Gauducheau ICO, Saint Herblain} &    &  \\
    
    \end{tabular}
    }

    \vspace{1.5cm}


    \begin{center}
    \textsc{Semestre 3 2024}
    \end{center}
    
\end{titlepage}
\let\cleardoublepage\clearpage


%%%%%%%%%%%%%%%%%%%%%%%%%%%%%%%%%%%%%%%%%%%%%%%%%%%%%%
%%%%%%%%%%%%%%%%%%%%%%% STYLE %%%%%%%%%%%%%%%%%%%%%%%%
%%%%%%%%%%%%%%%%%%%%%%%%%%%%%%%%%%%%%%%%%%%%%%%%%%%%%%

\onehalfspacing

%Style  du corps
\pagestyle{fancy}
	\renewcommand\headrulewidth{0.5pt}
	\renewcommand\footrulewidth{0.5pt}
	\fancyfoot[L]{\textsc{A. Rintaud}}
	\fancyfoot[C]{\textsc{ICO Nantes}}
	\fancyfoot[R]{\thepage}

\tableofcontents
\clearpage
\chapter{Introduction}

La radiothérapie externe utilise, de manière prépondérante, les faisceaux de photons de haute énergie afin de traiter des cellules cancéreuses tout en épargnant le plus possible les tissus sains. Dans cette optique, la mesure précise des caractéristiques dosimétriques, tels que les rendements en profondeur, les profils de dose, les facteurs d'ouverture du collimateur (FOC) ainsi que les incertitudes associées de l'accélérateur utilisé sont nécessaires en vue de fournir ces caractéristiques au TPS (\textit{Treatment Planning System}) utilisé en clinique pour le calcul des plans de traitement des patients.

Ce rapport traitera des faisceaux de photons utilisés en radiothérapie externe. Sera étudié l'impact de certains paramètres d'acquisition sur la distribution de dose (dose relative), ainsi que des paramètres d'influence (énergie du faisceau, tailles de champ, DSP, détecteurs etc).

\chapter{Matériels et méthodes}
\section{Grandeurs dosimétriques}
\subsection{Rendement en profondeur}

Le rendement en profondeur ($RDT$) permet de connaître l'évolution de la dose dans le milieu de référence en fonction de la profondeur $z$ du point de mesure, le long de l'axe du faisceau. Il est donné par la formule suivante :

\begin{equation}
    RDT(z, A, E) = \dfrac{D_z}{D_{max}} \times 100 \; (\%)
    \label{eq_rdt}
\end{equation}

Avec :

\begin{itemize}
    \item[$\bullet$] $z$ la profondeur du détecteur dans le fantôme d'eau
    \item[$\bullet$] $A$ la taille de champ au niveau de la surface
    \item[$\bullet$] $DSP$ la distance source-peau
    \item[$\bullet$] $E$ l'énergie du faisceau
    \item[$\bullet$] $D_z$ la dose à la profondeur $z$
    \item[$\bullet$] $D_{max}$ la dose maximale
\end{itemize}

La figure \ref*{fig_rdt} indique le montage pour l'acquisition du rendement en profondeur.

\begin{figure}[h]
  \centering
  \includegraphics[scale=0.3]{figures/rdt_schema.png}
  \caption{Géométrie pour la mesure du rendement en profondeur}
  \label{fig_rdt}
\end{figure}

Les rendements en profondeur possèdent trois zones caractéristiques :

\begin{itemize}
  \item[$\bullet$] zone de build-up (de la surface du fantôme jusqu'au point du dépôt de dose maximum) : cette zone où la dose relative est croissante jusqu'au maximum $z_{max}$ est la zone où les électrons secondaires sont mis en mouvement, l'équilibbre électronique n'est donc pas atteint et la dose est inférieure au kerma
  \item[$\bullet$] profondeur du maximum : profondeur à laquelle la dose déposée dans le milieu est maximale puisque l'équilibre électronique atteint
  \item[$\bullet$] zone de décroissance : zone où la dose diminue avec la profondeur engendrée par l'atténuation du faisceau dans le milieu
\end{itemize}

Dans la suite de ce rapport sera discuté de la dose à l'entrée du fantôme, la profondeur du dépôt de dose maximum ainsi que de la dose en profondeur.

% \subsection{Rapport tissu-fantôme}

% Le rapport tissu-fantôme (TPR) représente la dose en fonction de la profondeur normalisé à la profondeur de référence. Contrairement au rendement en profondeur, la distance source-détecteur est fixe comme le montre la figure \ref*{fig_tpr}. Le TPR s'exprime à l'aide de la formule \ref{eq_tpr} :

% \begin{equation}
%     TPR(z, A, DSP, E) = \dfrac{D_z}{D_{ref}} \times 100 \; (\%)
%     \label{eq_tpr}
% \end{equation}

% Avec $D_{ref}$ la dose à la profondeur de référence $z_{ref}$.

% \begin{figure}[h] 
%   \centering
%   \includegraphics[scale=0.3]{figures/tpr_schema.png}
%   \caption{Géométrie pour la mesure du TPR}
%   \label{fig_tpr}
% \end{figure}

% \subsection{Rapport tissu-maximum}

% Le RTM est un cas particulier du TPR puisqu'il est normalisé par rapport à la dose du maximum $z_{max}$. Il est défini par la fomule qui suit :

% \begin{equation}
%     RTM(z, A, E) = \dfrac{D_z}{D_{max}}
%     \label{eq_rtm}
% \end{equation}

% Comme la distance entre la source et le détecteur est constante, la mesure du RTM est difficile à réaliser. La formule \ref{eq_calcul_rtm} permet de calculer le RTM à partir du rendement en profondeur \cite{podgorsak2005radiation} :

% \begin{equation}
%     RTM(z, c, E) \approx RDT(z, c, DSP, E) \times \left(\dfrac{DSP + z}{DSP + z_{max}}\right)^2
%     \label{eq_calcul_rtm}
% \end{equation}

\subsection{Profils de dose}

Les profils de dose permettent d'évaluer la variation de la dose suivant les axes perpendiculaires à l'axe du faisceau qui sont nommés \textit{inline} et \textit{crossline}. Le profil \textit{inline} est orienté en gauche-droite et le \textit{crossline} est quant à lui orienté en tête-pieds. Les profils sont composés de trois régions principales (voir figure \ref*{fig_regions_profil}) pour des faisceaux avec un filtre égalisateur et sont normalisés par rapport à l'axe du faisceau :

\begin{itemize}
  \item[$\bullet$] la zone centrale, ou zone homogène définie par les points au dessus de l'isodose 80\% (voir Annexe \ref{partie_annexe} \cite{ansm2007})
  \item[$\bullet$] les régions de pénombre qui est la distance entre l'isodose 20\% et celle à 80\% de chaque côté du profil
  \item[$\bullet$] les zones hors champ qui correspond aux zones où les isodoses sont inférieures à 20\%
\end{itemize}

Pour caractériser les profils, plusieurs métriques sont à notre disposition :

\begin{itemize}
  \item[$\bullet$] L'homogénéité permet d'évaluer la planéité de la zone centrale du profil. Cette métrique se calcule comme suit, d'après le protocole choisi en clinique dans le logiciel MyQA :
  
  \begin{equation}
    H = \dfrac{D_{max} - D_{min}}{D_{max} + D_{min}} \times 100 \; (\%)
    \label{eq_homogeneite}
  \end{equation}
  
  \item[$\bullet$] La symétrie est définie de la manière suivante :
  
  \begin{equation}
    S = \max\left(\dfrac{D(-x)}{D(+x)}; \dfrac{D(+x)}{D(-x)}\right) \times 100 \; (\%)
    \label{eq_symetrie}
  \end{equation}
  Avec $x$ la distance à l'axe central.
  
  \item[$\bullet$] La pénombre est la distance séparant le point à 20\% de la dose à l'axe et celui à 80\%. Il y a donc une valeur pour chacun des côtés du faisceau.  
\end{itemize}

\begin{figure}[h]
  \centering
  \includegraphics[scale=0.8]{figures/regions_profil.PNG}
  \caption{Différentes régions d'un profil de dose}
  \label{fig_regions_profil}
\end{figure}

\subsection{FOC}

Le facteur d'ouverture du collimateur (FOC) est le rapport de la dose à une taille de champ $A$ quelconque et de la dose au champ de référence $A_{ref}$, qui sera de 1à$\times$10 cm$^2$ dans cette étude \footnote{La DSP et la profondeur de mesure $z$ sont fixes} (voir figure \ref*{fig_foc}):

\begin{equation}
    FOC(A, E, z) = \dfrac{D(A, E, z)}{D(A_{ref}, E, z)} \times 100 \; (\%)
    \label{eq_foc}
\end{equation}

Avec :

\begin{itemize}
    \item[$\bullet$] $D(A, E, z)$ la dose à la taille de champ $A$, l'énergie $E$ et à la profondeur $z$
    \item[$\bullet$] $D(A_{ref}, E, z)$ la dose à la taille de champ de référence $A_{ref}$, à l'énergie $E$ et à la profondeur $z$
\end{itemize}

\begin{figure}[h]
  \centering
  \includegraphics[scale=0.55]{figures/foc_schema.png}
  \caption{Géométrie de mesure des FOC \cite{mayles2007handbook}}
  \label{fig_foc}
\end{figure}

Comme l'indique la formule \ref*{eq_foc}, les FOC sont normalisés par rapport à la mesure avec un champ de 10$\times$10 cm$^2$. De plus, la profondeur et la DSP sont respectivement 10 cm et 90 cm lors des mesures de FOC pour cette étude.

\section{Matériels}

Les mesures de dose relative ont été réalisées à l'aide du logiciel MyQA (version 9.0.25.0) du constructeur IBA, associé à la cuve à eau BluePhantom 2 du même constructeur. Cette cuve à eau ainsi que MyQA permettent de piloter les déplacements d'un détecteur et de le déplacer dans les trois directions de l'espace. MyQA permet de créer également ds queues de mesures et d'analyser les mesures acquises. Le tableau \ref*{table_matos} regroupe l'ensemble du matériel utilisé pour ces mesures. Nous pouvons voir que dans le tableau \ref*{table_matos}, il y a deux chambres CC13 dites de champ et de référence. Celle dite de champ est pour l'acquisition du signal dans le fantôme d'eau, et la chambre dite de référence permet de s'affranchir des fluctuations du débit du faisceau. Concernant l'accélérateur, nous avons utilisé le Clinac iX2300 du constructeur Varian qui met a disposition deux faisceaux de photons de 6 MV et de 23 MV.

\begin{table}[h]
  \centering
  \begin{tabular}{>{\centering\arraybackslash}m{4cm}>{\centering\arraybackslash}m{3cm}>{\centering\arraybackslash}m{2.5cm}>{\centering\arraybackslash}m{2.5cm}>{\centering\arraybackslash}m{2cm}}
    \toprule
    \textbf{Matériel} & \textbf{Volume sensible (cm}$\mathbf{^3}$\textbf{)} & \textbf{Matériau} & \textbf{Constructeur} & \textbf{N}$\mathbf{^{\circ}}$\textbf{ de série}\\
    \toprule
    Chambre dite de référence CC13 & 0,13 & Air & IBA & 3922 \\
    Chambre dite de champ CC13 & 0,13 & Air & IBA & 3923 \\
    Chambre Farmer 30013 & 0,6 & Air & PTW & 011924 \\
    Chambre Semiflex 31010 & 0,125 & Air & PTW & 008214 \\
    Chambre Pinpoint 31014 & 0,04 & Air & PTW & 00787 \\
    Diode SRS 60018 & 0,03 & Silicium & PTW & 000186 \\
    MicroDiamant 60019 & 0,004 $\times$ 10$^{-3}$ & Diamant & PTW & 122271 \\
    Electromètre Unidos & / & / & PTW & 20505 \\
    Cuve à eau Blue Phantom 2 & / & / & IBA & 8173 \\
    Clinac iX 2300 (Clinac 2) & / & / & Varian & H294581 \\
    \bottomrule
  \end{tabular}
  \caption{Récapitulatif du matériel utilisé lors des mesures (acquisition continue)}
  \label{table_matos}
\end{table}

\newpage
\section{Méthodes}
\subsection{Mise en place de la cuve}

Cette partie traitera de l'installation de la cuve à eau (représentée sur la figure \ref*{fig_cuve}) pour les acquisitions de profils de dose et de rendements en profondeur, dont les étapes sont présentées par la suite :

\begin{itemize}
  \item[$\bullet$] Bras à 0$^{\circ}$ physique à l'aide d'un niveau 
  \item[$\bullet$] Alignement de la cuve avec le croisillon de l'accélérateur à l'aide d'un capuchon que l'on place sur le détecteur
  \item[$\bullet$] Remplissage de la cuve à eau
  \item[$\bullet$] Installation de la chambre de référence dans le tube prévu à cet effet
  \item[$\bullet$] Installation de la chambre de champ sur le socle permettant son déplacement dans les trois directions de l'espace
  \item[$\bullet$] Une fois que la cuve est remplie d'eau, aligner les rails en faisant se déplacer la chambre de champ munie de son capuchon\footnote{Le capuchon marqué d'une croix sur sa face perpendiculaire à la surface de l'eau ainsi que sur celle parrallèle à l'eau pour pouvoir l'aligner avec le croisillon} pour que celle-ci soit à la surface de l'eau sur toute la surface d'entrée du fantôme d'eau
  \item[$\bullet$] Enregistrer la position de l'isocentre en ayant au préalable aligné la chambre de champ avec le croisillon et au niveau de la surface de l'eau. 
  \item[$\bullet$] Sur le logiciel, bien vérifier l'orientation de la cuve pour que le repère de la cuve soit le même que celui du logiciel
  \item[$\bullet$] Placer la cuve à la DSP souhaitée
\end{itemize}

\begin{figure}[h]
  \centering
  \includegraphics[scale=0.4]{figures/cuve.png}
  \caption{Cuve à eau Blue Phantom 2 du constructeur IBA utilisée lors des mesures}
  \label{fig_cuve}
\end{figure}

L'instrallation de la cuve est une étape primordiale pour des mesures les plus précises possible puisque celles-ci sont importées dans le TPS qui permet la modélisation de l'accélérateur.

\newpage
\subsection{Mesures réalisées}

Pour chacune des grandeurs dosimétriques étudiées, nous allons réaliser plusieurs mesures en faisant varier plusieurs paramètres d'influences et d'aquisition comme le montre la figure \ref*{table_mesures}.

\begin{table}[h!]
  \centering
  \begin{tabular}{|>{\centering\arraybackslash}m{2.5cm}|>{\centering\arraybackslash}m{4cm}|>{\centering\arraybackslash}m{2.5cm}|>{\centering\arraybackslash}m{4cm}|}
    \hline
    \textbf{Grandeur dosimétrique} & \textbf{Paramètre d'influence} & \textbf{Mesures de référence} & \textbf{Comparaison}                    \\ \hline
    \multicolumn{1}{|c|}{\multirow{5}{*}{\textbf{Rendements}}} & \textbf{Taille de champ (cm}$\mathbf{^2}$\textbf{)} & 10x10 (machoires) & 3x3, 6x6, 20x20    (machoires)                \\
    \multicolumn{1}{|c|}{}                              & \textbf{DSP (cm)}             & 100                & 85, 110                                 \\
    \multicolumn{1}{|c|}{}                              & \textbf{Energie (MV)}         & 6                  & 23                                      \\
    \multicolumn{1}{|c|}{}                              & \textbf{Détecteur}            & CC13               & Pinpoint, MicroDiamant, Diode, Semiflex \\
    \multicolumn{1}{|c|}{}                              & \textbf{Chambre de référence} & Dans le champ      & Hors champ                              \\ \hline
    \multicolumn{1}{|c|}{\multirow{7}{*}{\textbf{Profils}}}    & \textbf{Champ (cm}$\mathbf{^2}$\textbf{)} & 10x10 & 3x3, 6x6, 8x8, 12x12, 15x15, 20x20 \\
    \multicolumn{1}{|c|}{}                              & \textbf{Profondeur (cm)}      & 10                 & 3, 20                                   \\
    \multicolumn{1}{|c|}{}                              & \textbf{Energie (MV)}         & 6                  & 23                                      \\
    \multicolumn{1}{|c|}{}                              & \textbf{DSP (cm)}             & 100                & 85, 110                                 \\
    \multicolumn{1}{|c|}{}                              & \textbf{Détecteur}            & CC13               & Pinpoint, MicroDiamant, Diode, Semiflex \\
    \multicolumn{1}{|c|}{}                              & \textbf{Orientation du profil}          & Crossline         & Inline                                  \\
    \multicolumn{1}{|c|}{}                              & \textbf{Mode d'acquisition}   & Continu & Step by step \\ \hline
    \multicolumn{1}{|c|}{\multirow{3}{*}{\textbf{FOC}}} & \textbf{Energie (MV)}         & 6                  & 23                                      \\
    \multicolumn{1}{|c|}{}                              & \textbf{Détecteur}            & CC13               & Farmer                                  \\
    \multicolumn{1}{|c|}{}                              & \textbf{DSP (cm)}             & 90                 & 80, 120                                 \\ \hline
  \end{tabular}
  \caption{Mesures réalisées dans le cadre de cette étude}
  \label{table_mesures}
\end{table}

Pour les acquisitions en modes continu et pas à pas, plusieurs combinaisons ont été étudiés. Le tableau \ref*{table_ss_presentation}

\begin{table}[h!]
  \centering
  \begin{tabular}{>{\centering\arraybackslash}m{2cm}>{\centering\arraybackslash}m{1.5cm}>{\centering\arraybackslash}m{2cm}>{\centering\arraybackslash}m{2cm}>{\centering\arraybackslash}m{3cm}}
    \toprule
    \textbf{Mode acquisition} & \textbf{Vitesse (cm/s)} & \textbf{Pas G-D (cm)} & \textbf{Pas centre (cm)} & \textbf{Temps d'intégration (s)} \\
    \toprule
    Continu & 0,3 & / & / & / \\
    Continu & 2,5 & / & / & / \\
    Pas à pas & / & 0,1 & 0,5 & 1 \\
    Pas à pas & / & 0,1 & 0,5 & 0,5 \\
    Pas à pas & / & 0,1 & 0,5 & 3 \\
    Pas à pas & / & 0,5 & 1 & 1 \\
    Pas à pas & / & 0,05 & 0,2 \\
    \bottomrule
  \end{tabular}
  \caption{Influence du mode d'acquisition sur les profils (résultats MyQA)}
  \label{table_ss_presentation}
\end{table}

Premièrement, les résultats concernant les rendements en profondeur seront présentés puis discutés. Il en sera de même pour les profils de dose et les facteurs d'ouverture du collimateur.

\newpage

\chapter{Résultats et discussions}
\section{Rendement en profondeur}
\subsection{Influence de l'énergie}

L'influence de l'énergie sur le rendement en profondeur pour un faisceau de photons est multiple. En effet, nous voyons tout d'abord sur la figure \ref*{fig_rdt_energie} et dans le tabelau \ref*{table_profils_energie} que la zone de mise en équilibre électronique est plus grande pour un faisceau de 23 MV que pour un faisceau de 6 MV. Cela s'explique par le fait que plus les photons primaires sont énergétiques, plus l'énergie transférée aux électrons secondaires sera importante. La zone de mise en équilibre électronique est donc plus grande (position du maximum de dépôt de dose plus profonde). De plus, nous observons que le rendement est plus important, à une profondeur donnée, pour des photons de 23 MV par rapport aux photons de 6 MV pour une profondeur supérieure au $z_{max}$. Cela est engendré par l'atténuation du milieu qui diminue pour une énergie qui augmente. Pour terminer, la dose à l'entrée est plus faible pour le faisceau le plus énergétique car la probabilité d'interaction est plus faible à la surface d'entrée. Il y a également moins de rayonnement diffusé et d'électrons de contamination dans la tête de l'accélérateur.

\begin{figure}[h!]
  \centering
  \includegraphics[scale=0.44]{../scripts/figures/dose_relative/rdt_X6_X23.png}
  \caption{Rendements en profondeur pour les différentes énergies de l'accélérateur utilisé pour une taille de champ de 10$\times$10 cm$^2$ et une DSP de 100 cm}
  \label{fig_rdt_energie}
\end{figure}

\begin{table}[h]
  \centering
  \begin{tabular}{cccc}
    \toprule
    \textbf{Energie (MV)} & $\mathbf{z_{max}}$ \textbf{(cm)} & $\mathbf{R_{50}}$ \textbf{(cm)} & \textbf{Dose surface (\%)} \\
    \toprule
    6 & 1,23 & 15,07 & 55,96 \\
    23 & 3,16 & 21,24 & 33,7 \\
    \bottomrule
  \end{tabular}
  \caption{Résultats de l'influence de l'énergie sur les rendements en profondeur pour les faisceaux de photons du Clinac 2}
  \label{table_rdt_energie}
\end{table}

\subsection{Influence de la DSP}

La figure \ref*{fig_rdt_dsp} nous montre l'influence de la DSP (distance source patient) sur le rendement en profondeur \footnote{La taille de champ à été adaptée pour qu'il y ait toujours un champ de 10x10 cm$^2$ à 10 cm de profondeur}. Nous pouvons voir que plus la DSP augmente, plus la dose relative est élevée.

L'influence de la DSP s'observe surtout en profondeur. La pente du rendement en profondeur est plus faible pour l'acquisition à 110 cm de DSP. Cela peut s'expliquer par l'effet de la loi d'inverse carrée de la distance, qui est plus important à distance plus faible, d'où une dose plus faible pour une DSP plus petite pour une même profondeur. 

Dans le tableau \ref*{table_profils_dsp}, nous voyons que la dose à la surface est plus importante pour une DSP qui diminue. Plus la DSP diminue, plus il y a des électrons de contamination et des photons diffusés de basse énergie émergeant de la tête de l'accélérateur puisque ces particules chargées légères diffusent dans l'air.

Concernant la profondeur du dépôt de dose, nous observons que le $z_{max}$ augmente avec la DSP (voir la figure \ref*{fig_rdt_DSP_zoom} zommée en Annexe). Cela s'explique par le fait que pour une DSP plus faible, le nombre d'électrons de contamination originaire de la tête de l'accélérateur sont plus nombreux (diffusion moindre des électrons dans l'air) ce qui implique une zone de build-up plus faible.

\begin{figure}[h]
  \centering
  \includegraphics[scale=0.45]{../scripts/figures/dose_relative/DSP/Rendement_DSP.png}
  \caption{Rendements en profondeur en faisant varier la DSP pour une taille de champ de 10$\times$10 cm$^2$ et une énergie de 6 MV}
  \label{fig_rdt_dsp}
\end{figure}

\begin{table}[h]
  \centering
  \begin{tabular}{cccc}
    \toprule
    \textbf{DSP (cm)} & $\mathbf{z_{max}}$ \textbf{(cm)} & $\mathbf{R_{50}}$ \textbf{(cm)} & \textbf{Dose surface (\%)} \\
    \toprule
    85 & 1,12 & 14,3 & 56,04 \\
    100 & 1,23 & 15,07 & 55,96 \\
    110 & 1,37 & 15,47 & 53.32 \\
    \bottomrule
  \end{tabular}
  \caption{Résultats de l'influence de la DSP sur les rendements en profondeur pour le faisceau de photons de 6 MV du Clinac 2}
  \label{table_rdt_dsp}
\end{table}

\subsection{Influence de la taille de champ}

La taille de champ influence la dose en profondeur, une fois que l'équilibre électronique est atteint. En effet, nous voyons sur la figure \ref*{fig_champs_profils} et le tableau \ref*{table_profils_champ} que plus la taille de champ augmente, plus la dose en profondeur augmente. Ce phénomène est expliqué par l'augmentation du volume diffusant. De plus, la dose à l'entrée augmente avec la taille de champ. Cela s'explique par une quantité plus importante de rayonnement diffusé dans la tête de l'accélérateur (surface du cône égalisateur irradiée plus importante et surface apparente des machoires plus grande).

Nous voyons que plus la taille de champ augmente, plus le $z_{max}$ diminue. Le fait d'augmenter la taille de champ augmente le nombre de photons diffusés et le nombre d'électrons de contamination. L'énergie moyenne du faisceau est donc plus faible, ce qui diminue la profondeur de mise en équilibre électronique. Il y a tout de même un souci sur la taille de champ de 20$\times$20 cm$^2$ sur la valeur du $z_{max}$. Il y a du y avoir un problème d'interpolation par le logiciel puisque sur la figure \ref*{fig_rdt_champ} nous voyons que le $z_{max}$ est la plus petite valeur des quatre courbes. Le résultat aurait surement été différent avec un détecteur muni d'un volume sensible de plus petite taille.

\begin{figure}[h!]
  \centering
  \includegraphics[scale=0.45]{../scripts/figures/dose_relative/rendement_champs.png}
  \caption{Rendements en profondeur pour différentes tailles champs. DSP de 100 cm et énergie de faisceau de 6 MV}
  \label{fig_rdt_champ}
\end{figure}

\begin{table}[h]
  \centering
  \begin{tabular}{cccc}
    \toprule
    \textbf{Champ (cm}$\mathbf{^2}$\textbf{)} & $\mathbf{z_{max}}$ \textbf{(cm)} & $\mathbf{R_{50}}$ \textbf{(cm)} & \textbf{Dose surface (\%)} \\
    \toprule
    3x3 & 1,48 & 12,92 & 48,11 \\
    6x6 & 1,36 & 13,94 &  50,66 \\
    10x10 & 1,23 & 15,07 & 55,96 \\
    20x20 & 1,36 & 16,7 & 62,63 \\
    \bottomrule
  \end{tabular}
  \caption{Résultats de l'influence de la DSP sur les rendements en profondeur pour le faisceau de photons de 6 MV du Clinac 2}
  \label{table_rdt_dsp}
\end{table}

\newpage
\subsection{Influence du détecteur}

La figure \ref*{fig_rdt_detecteurs} et le tableau \ref*{table_rdt_detecteurs} montrent l'influence du détecteur choisi sur le rendement en profondeur. Nous pouvons voir que la diode mesure une plus grande dose à la surface. Cela est dû au $Z$ effectif du détecteur qui est plus important que les autres détecteurs ($Z$ = 14), ce qui favorise l'effet photoélectrique à basse énergie donc un signal plus important. La diode SRS utilisée n'est donc pas blindée. Concernant les chambres d'ionisation Semiflex et CC13, leur volume sensible est très proche (0,13 cm$^3$ pour la CC13 et 0,125 cm$^3$ pour la Semiflex) ce qui explique leurs très fortes similarités (courbes quasiment superposées sur la figure \ref*{fig_rdt_detecteurs}). Nous observons que le signal de la Pinpoint et du MicroDiamant sont bruités. Leur volume sensible étant très petits (0,015 cm$^3$ et 0,004 mm$^3$ respectivement) la statistique de comptage est plus faible sur ces deux détecteurs puisque la vitesse d'acquisition est identique entre chacun des détecteurs. Il aurait donc fallu diminuer la vitesse d'acquisition pour ces détecteurs pour obtenir une meilleure statistique de comptage.

\begin{figure}[h]
  \centering
  \includegraphics[scale=0.45]{../scripts/figures/dose_relative/test/Rendement.png}
  \caption{Rendement en profondeur pour différents détecteurs. Champ de 10$\times$10 cm$^2$ et DSP de 100 cm}
  \label{fig_rdt_detecteurs}
\end{figure}

\begin{table}[h]
  \centering
  \begin{tabular}{cccc}
    \toprule
    \textbf{Détecteur} & $\mathbf{z_{max}}$ \textbf{(cm)} & $\mathbf{R_{50}}$ \textbf{(cm)} & \textbf{Dose surface (\%)} \\
    \toprule
    CC13 & 1,35 & 14,98 & 55,96 \\
    Diode & 0,7 & 14,55 & 90,72 \\
    MicroDiamant & 1,57 & 14,72 & 55,52 \\
    Pinpoint & 1,46 & 15,65 & 49,86 \\
    Semiflex & 1,29 & 14,96 & 55,68 \\
    \bottomrule
  \end{tabular}
  \caption{Résultats de l'influence du détecteur sur les rendements en profondeur pour le faisceau de photons de 6 MV du Clinac 2}
  \label{table_rdt_detecteurs}
\end{table}

\newpage
\subsection{Influence de la chambre de référence}

\begin{figure}[h!]
  \centering
  \includegraphics[scale=0.4]{../scripts/figures/dose_relative/chambre_ref.png}
  \caption{Rendements en profondeur avec et sans chambre de référence pour une taille de champ de 10$\times$10 cm$^2$ et une DSP de 100 cm}
  \label{fig_sanss_chambre_ref}
\end{figure}

Pour s'affranchir de la fluctuation du débit de dose du faisceau lors de la mesure de dose relative, une chambre dite de référence est placée dans l'air, dans un coin du champ d'irradiation pour ne pas perturber la mesure avec la chambre dans la cuve à eau. Nous avons donc réalisé une mesure de rendement en profondeur pour observer ce rendement sans que la chambre de référence soit placée dans le champ. La figure \ref*{fig_sanss_chambre_ref} permet de voir l'importance de la présence de la chambre de référence. En effet, nous voyons que le signal est très bruité et ne permet donc pas d'analyser correctement les résultats.

Cependant, une erreur a été faite lors de cette acquisition puisque la chambre de référence n'a pas été déconnectée du logiciel. Ceci implique que la chambre continuait de mesurer du signal hors champ, ce qui amplifie énormément les fluctuations du signal réceptionné par la chambre de champ ainsi que le bruit statistique.

\section{Profils de dose}
\subsection{Influence de l'énergie}

La pénombre augmente avec l'énergie, comme nous pouvons le voir sur la figure \ref*{fig_profils_energie} et le tableau \ref*{table_profils_energie}, car la transmission à travers les machoires est plus importante et que, pour des photons de plus haute énergie, l'énergie des particules secondaires est plus importante également, ce qui augmente la pénombre. De plus, nous voyons que les deux profils n'ont pas le même épaulement. Cela est dû au cône égalisateur qui est différent entre le faisceau de 6 MV et celui de 23 MV. Le cône atténue moins sur les côtés par rapport au centre car il y a moins de matière à traverser (faisceau plus dur au centre). Le faisceau étant moins dur sur les côtés, l'énergie moyenne est plus faible qu'au centre, ce qui favorise l'effet photoélectrique et donc une surestimation de la dose sur les côtés.

\begin{figure}[h]
  \centering
  \includegraphics[scale=0.45]{../scripts/figures/dose_relative/profils/profils_energie.png}
  \caption{Profils de dose pour différentes énergies 
  . Taille de champ de 10$\times$10 cm$^2$ et DSP de 100 cm}
  \label{fig_profils_energie}
\end{figure}

\begin{table}[h]
  \centering
  \begin{tabular}{>{\centering\arraybackslash}m{1.5cm}>{\centering\arraybackslash}m{2cm}>{\centering\arraybackslash}m{2cm}>{\centering\arraybackslash}m{2.5cm}>{\centering\arraybackslash}m{2.3cm}>{\centering\arraybackslash}m{2.5cm}}
    \toprule
    \textbf{Energie (MV)} & \textbf{Symétrie (\%)} & \textbf{Homogénéité (\%)} & \textbf{Centre du champ (cm)} & \textbf{Pénombre (G-D) (cm)} & \textbf{Taille de champ (cm)} \\
    \toprule
    6 & 101,38 & 2,55 & -0,01 & 0,76 - 0,755 & 11,13 \\
    23 & 101,15 & 2,5 & 0,03 & 0,87 - 0,86 & 11,15 \\
    \bottomrule
  \end{tabular}
  \caption{Influence de l'énergie sur les profils (résultats MyQA)}
  \label{table_profils_energie}
\end{table}

\newpage
\subsection{Influence de la DSP}

Ayant adapté la taille de champ pour chacune des DSP dans le but d'avoir la même taille de champ de 10$\times$10 cm$^2$ à la profondeur de mesure, la figure \ref*{fig_profils_DSP} et le tableau \ref*{table_profils_dsp} montrent très peu de différences. Cela s'explique par le fait que l'air est un milieu très peu diffusant et atténuant pour des faisceaux de photons de haute énergie. Nous pouvons donc considérer que la mesure de profils à profondeur constante en faisant varier la DSP n'influe que très peu sur la mesure, pour une profondeur donnée.

Nous observons que pour la DSP la plus faible, le signal est plus bruité que les deux autres acquisitions. La chambre de référence était potentiellement mal positionnée vue que nous avons adapté la taille de champ entre chaque acquisition (oublie de repositionner la chambre de référence dans le coin du champ).

\begin{figure}[h]
  \centering
  \includegraphics[scale=0.45]{../scripts/figures/dose_relative/DSP/Crossline_DSP.png}
  \caption{Profils de dose pour différentes DSP. orientation \textit{crossline}, profondeur de 10 cm et taille de champ de 10$\times$10 cm$^2$}
  \label{fig_profils_DSP}
\end{figure}



\begin{table}[h]
  \centering
  \begin{tabular}{>{\centering\arraybackslash}m{1.5cm}>{\centering\arraybackslash}m{2cm}>{\centering\arraybackslash}m{2cm}>{\centering\arraybackslash}m{2.5cm}>{\centering\arraybackslash}m{2.3cm}>{\centering\arraybackslash}m{2.5cm}}
    \toprule
    \textbf{DSP (cm)} & \textbf{Symétrie (\%)} & \textbf{Homogénéité (\%)} & \textbf{Centre du champ (cm)} & \textbf{Pénombre (G-D) (cm)} & \textbf{Taille de champ (cm)} \\
    \toprule
    85 & 101,28 & 2,35 & 0 & 0,67 - 0,67 & 11,08 \\
    100 & 100,43 & 2,37 & 0,01 & 0,68 - 0,68 & 11,07 \\
    110 & 101,81 & 2,39 & -0,04 & 0,69 - 0,69 & 11,12 \\
    \bottomrule
  \end{tabular}
  \caption{Influence de la DSP sur les profils (résultats MyQA)}
  \label{table_profils_dsp}
\end{table}

\newpage
\subsection{Influence de la taille de champ}

Le tableau \ref*{table_profils_champ} montre que la pénombre augmente avec la taille de champ. Ce phénomène vient du fait que la surface apparente du cône égalisateur est plus élevée à grande ouverture par rapport à un petit champ, ce qui augmente le rayonnement diffusé dans la tête de l'accélérateur.

De plus, le fait d'augmenter la taille de champ augmente la surface apparente du collimateur, ce qui implique une augmentation du diffusé et donc de la pénombre. Le volume diffusant dans le fantôme augmente avec la taille de champ, ce qui implique une plus grande contribution du rayonnnement diffusé dans la cuve d'eau, la pénombre augmente également avec le phénomène physique. Nous pouvons voir également sur le tableau \ref*{table_profils_champ} que l'homogénéité se dégrade lorsque la taille de champ diminue. Cette grandeur est mesurée sur la partie plane du profil. Or, pour des petits champs, le profil n'est plus vraiment plat mais se rapproche de la forme d'une gaussienne (voir figure \ref*{fig_champs_profils}), il est donc comprehensible que le résultat de l'homogénéité soit affecté pour de telles tailles de champ. 


\begin{figure}[h]
  \centering
  \includegraphics[scale=0.4]{../scripts/figures/dose_relative/profils_champs/comp_champs.png}
  \caption{Profils de dose pour différentes tailles de champs. Orientation \textit{crossline}, profondeur de 10 cm et DSP de 100 cm}
  \label{fig_champs_profils}
\end{figure}

Nous pouvons voir que le champ de 20$\times$20 cm$^2$ est un peu plus épaulé que les champs de plus petites dimensions. Comme discuté précédemment dans la partie \ref*{partie_rdt_energie}, le faisceau est moins dur sur les bords puisqu'il y a moins de d'épaisseur de filtre à traverser, ce qui va déposer plus de dose sur les bords par rapport à la dose à l'axe pour une même profondeur.

\begin{table}[h]
  \centering
  \begin{tabular}{>{\centering\arraybackslash}m{2cm}>{\centering\arraybackslash}m{1.5cm}>{\centering\arraybackslash}m{2cm}>{\centering\arraybackslash}m{2.3cm}>{\centering\arraybackslash}m{2.1cm}>{\centering\arraybackslash}m{2.3cm}}
    \toprule
    \textbf{Champ théorique (cm}$\mathbf{^2}$ \textbf{)} & \textbf{Symétrie (\%)} & \textbf{Homogénéité (\%)} & \textbf{Centre du champ (cm)} & \textbf{Pénombre (G-D) (cm)} & \textbf{Taille de champ (cm)} \\
    \toprule
    3x3 & 100,14 & 7,77 & 0 & 0,56 - 0,55 & 3,36 \\
    6x6 & 100,59 & 2,91 & 0,01 & 0,62 - 0,62 & 6,63 \\
    8x8 & 100,97 & 2,49 & -0,02 & 0,67 - 0,67 & 8,85 \\
    10x10 & 100,43 & 2,37 & 0,01 & 0,68 - 0,68 & 11,07 \\
    12x12 & 100,77 & 2,02 & -0,03 & 0,71 - 0,72 & 13,24 \\
    15x15 & 101,1 & 2,02 & 0,02 & 0,76 - 0,75 & 16,56 \\
    20x20 & 100,82 & 1,9 & -0,02 & 0,83 - 0,84 & 22,08 \\
    \bottomrule
  \end{tabular}
  \caption{Influence de la taille de champ sur les profils (résultats MyQA)}
  \label{table_profils_champ}
\end{table}

\newpage
\subsection{Influence du détecteur}

La figure \ref*{fig_profils_detecteurs} et le tableau \ref{table_profils_detecteurs} nous montrent que la pénombre est très dépendante du détecteur. En effet, le choix du détecteur est crucial lorsqu'il y a de forts gradients de dose puisqu'en fonction du volume sensible du détecteur, les pentes des pénombres vont être modifiées. Nous pouvons voir qu'avec un détecteur possédant un volume sensible très petit (comme le MicroDiamant et la diode SRS) que la pénombre est petite. Si le volume sensible du détecteur est trop important, le signal reccueilli n'est pas représentatif de la réalité puisque le signal va être moyenné sur une trop grande distance. Cependant, nous remarquons que pour la chambre d'ionisation Pinpoint, qui possède un volume sensible très petit de 0,015 cm$^3$, la dose est surrestimée au niveau des queues de distribution. Le bruit engendré par l'irradiation du manche de la chambre n'est plus négligeable par rapport au faible signal reccueilli par le petit volume sensible. Cela a pour effet d'augmenter la pénombre, ce qui est observable dans le tableau \ref{table_profils_detecteurs} (0,65 cm de pénombre pour la Pinpoint contre 0,39 en moyenne pour le MicroDiamant).

Nous voyons sur la figure \ref*{fig_profils_detecteurs} que certaines acquisitions sont bruitées (PinPoint, MicroDiamant et diode). La vitesse d'acquisition est sans doute trop importante, il aurait fallu ajuster la valeur lors des mesures.

Concernant le signal hors champ de la diode, ce détecteur ayant un $Z$ effectif étant plus élevé que celui des autres détecteurs, le rayonnement diffusé hors champ est plus facilement détecté avec la diode (effet photoélectrique favorisé). De plus, cette diode n'étant pas blindée, les photons de basse énergie ne sont pas filtrés par le blindage, ce qui augmente le signal reçu.

Pour le détecteur MicroDiamant nous voyons que sa réponse hors champ est plus faible que pour le reste des détecteurs. Cela peut s'expliquer par le très faible volume sensible (0,004 mm$^3$). Le signal détecté est bien plus faible que pour une chambre d'ionisation ou une diode par exemple.

\begin{figure}[h]
  \centering
  \includegraphics[scale=0.45]{../scripts/figures/dose_relative/test/profils_10x10.png}
  \caption{Profils de dose pour différents détecteurs. Taille de champ de 10$\times$10 cm, orientation \textit{crossline}, profondeur de 10 cm et DSP de 100 cm}
  \label{fig_profils_detecteurs}
\end{figure}

\begin{table}[h]
  \centering
  \begin{tabular}{>{\centering\arraybackslash}m{1.7cm}>{\centering\arraybackslash}m{2cm}>{\centering\arraybackslash}m{2cm}>{\centering\arraybackslash}m{2.5cm}>{\centering\arraybackslash}m{2.2cm}>{\centering\arraybackslash}m{3cm}}
    \toprule
    \textbf{Détecteur} & \textbf{Symétrie (\%)} & \textbf{Homogénéité (\%)} & \textbf{Centre du champ (cm)} & \textbf{Pénombre G-D (cm)} & \textbf{Taille de champ (cm)} \\
    \toprule
    CC13 & 100,56 & 2,25 & -0,05 & 0,68 - 0,68 & 11,07 \\
    Semiflex & 100,9 & 2,53 & -0,04 & 0,67 - 0,67 & 11,1 \\
    Diode & 101,01 & 2,64 & -0,02 & 0,45 - 0,44 & 11,11 \\
    MicroDiamant & 101,72 & 2,91 & -0,04 & 0,38 - 0,4 & 11,04 \\
    Pinpoint & 101,99 & 2,68 & -0,02 & 0,65 - 0,65 & 11,07 \\
    \bottomrule
  \end{tabular}
  \caption{Influence du détecteur sur les profils (résultats MyQA)}
  \label{table_profils_detecteurs}
\end{table}

\newpage
\subsection{Influence de l'orientation du profil}

Lors de l'acquisition des profils, nous pouvons choisir l'axe selon lequel le profil sera enregistré. Les résultats des deux profils sont donnés dans la tableau \ref*{table_profils_orientation} et sur la figure \ref*{fig_orientation_profil}. Les résultats montrent que la pénombre est modifiée entre ces deux acquistions. En effet, la pénombre est plus importante en \textit{inline} (tête-pieds) par rapport à l'orientation \textit{crossline} (droite-gauche). C'est la prosition des machoîres qui sont responsables de cette différence car les deux paires de machoîres ne sont pas sur le même plan. Celles qui définissent le champ en \textit{crossline} (nommées Y) sont en-dessous de celles qui le définissent en \textit{inline} (nommées X). Il s'agit donc d'un phénomène géométrique.

De plus, le tableau \ref*{table_profils_orientation} montre une légère différence de symétrie entre les deux orientations de profils. Ceci peut s'expliquer par le fait que l'orientation du faisceau est géré par deux systèmes de bobines. Un système pour l'orientation tête-pieds et l'autre pour le droite-gauche, ce qui peut impacter la symétrie du faisceau.




l'irradiation du détecteur n'est pas uniforme en \textit{inline} puisque la chambre est parallèle au profil acquis. Ce qui n'est pas le cas en \textit{crossline}.

\begin{figure}[h]
  \centering
  \includegraphics[scale=0.44]{../scripts/figures/dose_relative/orientation_profil.png}
  \caption{Profils de dose pour différentes orientations du profil. Taille de champ de 10$\times$10 cm, orientation \textit{crossline}, profondeur de 10 cm et DSP de 100 cm}
  \label{fig_orientation_profil}
\end{figure}

\begin{table}[h]
  \centering
  \begin{tabular}{>{\centering\arraybackslash}m{1.7cm}>{\centering\arraybackslash}m{2cm}>{\centering\arraybackslash}m{2cm}>{\centering\arraybackslash}m{2.5cm}>{\centering\arraybackslash}m{2.5cm}>{\centering\arraybackslash}m{3cm}}
    \toprule
    \textbf{Orientation} & \textbf{Symétrie (\%)} & \textbf{Homogénéité (\%)} & \textbf{Centre du champ (cm)} & \textbf{Pénombre (G-D) (cm)} & \textbf{Taille de champ (cm)} \\
    \toprule
    Inline & 101,38 & 2,55 & -0,01 & 0,76 - 0,75 & 11,13 \\
    Crossline & 100,43 & 2,37 & 0,01 & 0,68 - 0,68 & 11,07 \\
    \bottomrule
  \end{tabular}
  \caption{Influence de l'orientation du profil (résultats MyQA)}
  \label{table_profils_orientation}
\end{table}

\newpage
\subsection{Influence du mode d'acquisition}

Pour observer l'impact du mode d'acquisition sur les profils de dose, nous avons réalisé plusieurs mesures, comme nous plouvons le voir dans le tableau \ref*{table_ss}.

% \begin{itemize}
%   \item[$\bullet$] Continu avec une vitesse de déplacement de chambre de 0,3 cm/s
%   \item[$\bullet$] Continu avec une vitesse de déplacement de chambre de 2,5 cm/s
%   \item[$\bullet$] Pas à pas avec une mesure tous les 0,1 cm sur les côtés et une mesure centrale tous les 0,5 cm avec un temps d'intéragtion d'une seconde pour chaque point
%   \item[$\bullet$] Pas à pas avec une mesure tous les 0,1 cm sur les côtés et une mesure centrale tous les 0,5 cm avec un temps d'intéragtion d'une demi seconde pour chaque point
%   \item[$\bullet$] Pas à pas avec une mesure tous les 0,1 cm sur les côtés et une mesure centrale tous les 0,5 cm avec un temps d'intéragtion de trois secondes pour chaque point
%   \item[$\bullet$] Pas à pas avec une mesure tous les 0,5 cm sur les côtés et une mesure centrale tous les centimètres avec un temps d'intéragtion d'une seconde pour chaque point
%   \item[$\bullet$] Pas à pas avec une mesure tous les 0,05 cm sur les côtés et une mesure centrale tous les 0,2 cm avec un temps d'intéragtion d'une seconde pour chaque point
% \end{itemize}

% \begin{figure}[h]
%   \centering
%   \begin{subfigure}{\textwidth}
%     \centering
%     \includegraphics[scale=0.35]{../scripts/figures/dose_relative/step_by_step/step_by_step_1.png}
%   \end{subfigure}
%   \vspace{0.5cm}
%   \begin{subfigure}{\textwidth}
%     \centering
%     \includegraphics[scale=0.36]{../scripts/figures/dose_relative/step_by_step/step_by_step_2.png}
%   \end{subfigure}
%   \caption{Influence du mode d'acquisition sur les profils de dose}
%   \label{fig_mode_acquisition}
% \end{figure}

\begin{table}[h]
  \centering
  \begin{tabular}{>{\centering\arraybackslash}m{2cm}>{\centering\arraybackslash}m{1cm}>{\centering\arraybackslash}m{2cm}>{\centering\arraybackslash}m{1cm}>{\centering\arraybackslash}m{1.5cm}>{\centering\arraybackslash}m{0.7cm}>{\centering\arraybackslash}m{0.7cm}>{\centering\arraybackslash}m{1.7cm}>{\centering\arraybackslash}m{1.7cm}}
    \toprule
    \textbf{Mode acquisition} & \textbf{Vitesse (cm/s)} & \textbf{Pas G-D (cm)} & \textbf{Pas centre (cm)} & \textbf{Temps d'intégration (s)} & \textbf{S (\%)} & \textbf{H (\%)} & \textbf{Taille de champ (cm)} & \textbf{Pénombre G-D (cm)}\\
    \toprule
    Continu & 0,3 & / & / & / & 2,36 & 100,86 & 11,08 & 0,67 - 0,68\\
    Continu & 2,5 & / & / & / & 2,51 & 100,83 & 11,08 & 0,69 - 0,69\\
    Steb by Step & / & 0,1 & 0,5 & 1 & 2,23 & 100,67 & 11,08 & 0,66 - 0,67\\
    Step by Step & / & 0,1 & 0,5 & 0,5 & 2,26 & 100,53 & 11,08 & 0,67 - 0,67\\
    Step by step & / & 0,1 & 0,5 & 3 & 2,27 & 100,66 & 11,08 & 0,68 - 0,67\\
    Step by Step & / & 0,5 & 1 & 1 & 2,41 & 100,78 & 11,11 & 0,80 - 0,85\\
    Step by step & / & 0,05 & 0,2 & 1 & 2,26 & 100,61 & 11,08 & 0,66 - 0,67 \\
    \bottomrule
  \end{tabular}
  \caption{Inlfuence du mode d'acquisition sur les profils (résultats MyQA)}
  \label{table_ss}
\end{table}

Premièrement, concernant le mode continu, nous voyons sur le tableau \ref*{table_ss} que la différence de pénombre entre les vitesses lentes et rapides est relativement faible (0,015 cm de différence en moyenne pour la chambre CC13). Choisir une vitesse intermédiaire est satisfaisant pour des mesures rapides et précises. De plus, nous pouvons voir que le mode pas à pas avec une mesure tous les 0,05 cm sur les côtés du profils et tous les 0,2 cm au centre permet d'obtenir la plus petite pénombre mais n'est pas significatif pour les modes intermédiaires. Cependant, choisir un pas d'intégration trop grand n'est pas non plus optimal car la pénombre mesurée est plus grande mais également non symétrique car le logiciel doit interpoler entre les points pour obtenir les abscisses des points à 20\% et à 80\% de la dose, ce qui fausse la mesure de la pénombre. Il faut cependant bien mettre en perspective que ces résultats sont liés uniquement à l'utilisation de la chambre CC13 qui possède un volume sensible relativement grand. Les mêmes mesures ne seront certainement pas les mêmes avec un détecteur du type MicroDiamant ou diode.

\section{Présence de la chambre de référence}

La figure \ref*{fig_perturbations_chambre_ref_profils} et le tableau \ref*{table_resultats_perturb_chambre_ref} nous montrent les résultats de l'influence de la position de la chambre de référence dans le champ d'irradiation. Comme nous pouvons le voir sur le graphe et le tableau, le fait que la chambre de référence soit placée dans le coin ou presque au centre du champ ne change rien sur les résultats de l'acquisition puisque les résultats fournis par le tableau \ref*{table_resultats_perturb_chambre_ref} nous montrent que la symétrie, l'homogénéité et la pénombre sont extêmement proches. Nous pouvons conclure que la position de la chambre de référence perturbe très peu la signal reccueilli par la chambre de mesure.

\begin{figure}[h!]
  \centering
  \includegraphics[scale=0.41]{../scripts/figures/dose_relative/perturbations_chambre/perturbations_X6.png}
  \caption{Profils de dose pour différentes prositions de la chambre de référence. Taille de champ de 1O$\times$10 cm$^2$, profondeur de 10 cm et DSP de 100 cm}
  \label{fig_perturbations_chambre_ref_profils}
\end{figure}

\begin{table}[h]
  \centering
  \begin{tabular}{>{\centering\arraybackslash}m{2.5cm}>{\centering\arraybackslash}m{1.1cm}>{\centering\arraybackslash}m{1.1cm}>{\centering\arraybackslash}m{2.5cm}>{\centering\arraybackslash}m{2cm}>{\centering\arraybackslash}m{2.5cm}}
    \toprule
    \textbf{Position chambre ref} & \textbf{S (\%)} & \textbf{H (\%)} & \textbf{Centre du champ (cm)} & \textbf{Pénombre G-D (cm)} & \textbf{Taille de champ (cm)} \\
    \toprule
    Coin & 100,43 & 2,37 & 0,01 & 0,68 - 0,68 & 11,07 \\
    Centre & 100,69 & 2,29 & 0 & 0,68 - 0,68 & 11,12 \\
    \bottomrule
  \end{tabular}
  \caption{Influence de la position de la chambre de référence sur les profils de dose (résultats MyQA)}
  \label{table_resultats_perturb_chambre_ref}
\end{table}

\section{FOC}

Le tableau \ref*{table_resultats_foc} et la figure \ref*{fig_foc} nous montre les FOC mesurés pour des champs allant de 3x3 cm$^2$ à 20x20 cm$^2$. Premièrement, nous voyons que pour des champs inférieurs à 10$\times$10 cm$^2$ les FOC sont plus importants pour le faisceau de 23 MV que pour celui de 6 MV pour un détecteur donné. Cela peut s'expliquer par le fait qu'à plus faible énergie la quantité de photons et d'électrons qui rétrodiffusent sur les machoires vers les chambres monitrices dans la tête de l'accélérateur est plus importante (pour de petits champs), la dose mesurée par les chambres monitrices sera plus importante que la réalité et la coupure du faisceau sera donc prématurée. A contrario, pour une grande taille de champ (supérieure à 10$\times$10 cm$^2$), ce phénomène devient de plus en plus négligeable et la quantité de photons diffusés dans le fantôme d'eau à plus faible énergie prend le dessus. Ceci explique pourquoi les FOC sont plus faible pour une énergie de 23 MV plutôt qu'à 6 MV pour les grandes tailles de champ (et inversement).

Lorsqu'on utilise une chambre d'ionisation avec un volume sensible élevé (chambre de type Farmer), les résultats des FOC sont quasiment identiques par rapport à la chambrer CC13, à l'exception du champ 3x3 cm$^2$ (les faisceaux sont considérés comme petits à partir de 4$\times$4 cm$^4$ environ). Cela est dû au trop grand volume sensible du détecteur par rapport aux dimensions du champ, la mesure est donc sous-estimée avec la chambre Farmer.

Enfin,  nous observons que le rôle de la DSP, pour la profondeur de mesure, n'est pas significatif. En effet, la taille de champ ayant été adaptée, nous pouvons considérer que l'interaction du faisceau avec l'air (entre la source et le fantôme) est négligeable.

\begin{figure}[h]
  \centering
  \includegraphics[scale=0.45]{../scripts/figures/dose_relative/FOC.png}
  \caption{Facteurs d'ouverture du collimateur (FOC)}
  \label{fig_foc}
\end{figure}

\begin{table}[t!]
  \centering
  \begin{tabular}{>{\centering\arraybackslash}m{1.5cm}>{\centering\arraybackslash}m{1.8cm}>{\centering\arraybackslash}m{1.8cm}>{\centering\arraybackslash}m{1.8cm}>{\centering\arraybackslash}m{1.8cm}>{\centering\arraybackslash}m{1.8cm}>{\centering\arraybackslash}m{1.8cm}}
  \toprule
  \textbf{Champ (cm}$\mathbf{^2}$\textbf{)} &
  \textbf{FOC ref (\%)} &
  \textbf{FOC X23 (\%)} &
  \textbf{FOC Farmer (\%)} &
  \textbf{FOC DSP 80 (\%)} &
  \textbf{FOC DSP 120 (\%)} \\ \toprule
  3x3      & 82,44 & 84,13 & 81,12 & 82,54 & 83,35 \\
  6x6      & 91,71 & 94,38  & 91,74 & 91,97 & 91,88 \\
  8x8      & 96,18 & 97,74 & 96,38 & 96,47  & 96,39 \\
  10x10 & 100      & 100      & 100      & 100      & 100      \\
  12x12 & 102,93 & 101,83 & 102,99 & 102,79 & 102,89 \\
  15x15    & 106,51 & 103,79 & 106,45 & 106,27 & 106,59 \\
  20x20 & 110,94 & 106,06 & 110,55 & 110,45 & 111,08 \\ \bottomrule
  \end{tabular}
  \caption{Résultats des mesures des FOC}
  \label{table_resultats_foc}
\end{table}

\clearpage
\chapter{Annexe}
\label{partie_annexe}

\section{Définition de la zone homogène d'un champ d'irradiation en régime photons}

La zone homogène d’un champ d’irradiation en régime photons, pour les champs carrés de dimensions supérieures ou égales à 10 cm$\times$10 cm, est définie par la surface contenant des points distants de l’axe d’au plus :
\begin{itemize}
  \item[$\bullet$] 80\% de la demi-longueur des axes médians du champ lumineux, pour des points situés sur les médianes
  \item[$\bullet$] 70\% de la demi-longueur des diagonales du champ lumineux, pour des points situés sur les diagonales, la dose étant normalisée à 100\% sur l’axe du faisceau.
\end{itemize}
Ces pourcentages sont ramenés à la valeur commune de 60\% pour les champs carrés de dimensions comprises entre 5 cm $\times$ 5 cm et 10 cm $\times$ 10 cm.

\section{Influence de la DSP sur le rendement en profondeur}

\begin{figure}[h!]
  \centering
  \includegraphics[scale=0.4]{../scripts/figures/dose_relative/DSP/Rendement_DSP_zoom.png}
  \caption{Rendements en profondeur en faisant varier la DSP pour une taille de champ de 10$\times$10 cm$^2$ et une énergie de 6 MV (figure zoomée)}
  \label{fig_rdt_DSP_zoom}
\end{figure}

\clearpage
\bibliography{biblio}
\addcontentsline{toc}{chapter}{Références}
\bibliographystyle{plain}
\nocite{*}

\end{document}